% 第九节：交比、对合与调和点列
\chapter{交比、对合与调和点列}

\section{子空间上的矢量：嵌入、限制与射影}

前几节讨论的矢量都取在三维矢量空间 $V$ 中，对应平面上的点与线。但平面几何中有一类问题只涉及
一条直线上的点，或只涉及过一个定点的直线，前者称为点列，后者称为线束。直线上的点集是一维的，
过定点的直线集也是一维的，因此描述它们最自然的位置是一个二维矢量空间，而不是三维的 $V$。

设平面上的点用 $V$ 中的点矢 $a$ 描述，并要求点在一条固定的直线 $l$ 上，即
\[
l\cdot a=l_ia^i=0 .
\]
满足这一条件的全体点矢
\[
S=\{a\in V\mid l_ia^i=0\}
\]
构成 $V$ 的一个二维子空间：$S$ 对加法与数乘封闭，且恰好是一个过原点的平面。于是可以
把 $S$ 当作一个独立的二维矢量空间来研究，只在必要时才把它与 $V$ 中的平面等同起来。也就是说，
我们只关心到同构的层面。

\begin{definition}[嵌入]
设 $S\subset V$ 是二维子空间。把 $S$ 中的元素看作 $V$ 的元素，所得到的包含映射
\[
\iota:S\hookrightarrow V
\]
称为从 $S$ 到 $V$ 的\emph{嵌入}。通过嵌入，$S$ 中的二维点矢被实现为 $V$ 中的三维点矢，
几何上相当于把直线上的点画到平面上。
\end{definition}

嵌入把点矢从二维提升到三维。反方向的操作作用在对偶矢量上，由对偶矢量的函数意义给出。

\begin{definition}[限制]
设 $S\subset V$ 是二维子空间，$f\in V^*$ 是 $V$ 上的线性函数。把 $f$ 的定义域限制为
$S$，得到 $S$ 上的线性函数
\[
f_S:S\to\R,\qquad f_S(a)=f(a)\quad(a\in S),
\]
称为 $f$ 在 $S$ 上的\emph{限制}。$f\mapsto f_S$ 给出线性映射 $V^*\to S^*$。
\end{definition}

一句话，限制映射就是限制了线矢作为函数的定义域。由对偶矢量的函数意义，这一操作是自然的：
两个在 $V$ 中不同的线矢，限制到 $S$ 上之后可能给出同一个 $S^*$ 中的元素，因此限制一般不可逆。

嵌入与限制地位对称，但都不够直观。相比之下较为直观的是从线矢到点矢的映射。

\begin{definition}[射影]
设 $S$ 是对应于直线 $l$ 的二维子空间，$m$ 是任意一条不与 $l$ 平行的直线（即
$m$ 与 $l$ 不共向，二者在平面内交于唯一一点）。令
\[
a=l\times m,
\]
即把线矢 $m$ 映射为 $l$ 与 $m$ 的交点。这一映射 $V^*\to S$ 称为\emph{射影}。
\end{definition}

取分量写法。设 $e^{jk}=\eps^{ijk}l_i$，则由叉乘定义
\[
a^k=(l\times m)^k=\eps^{kij}l_im_j=m_j\eps^{jik}l_i=m_je^{jk},
\]
其中用到 $\eps^{jik}=-\eps^{ijk}$ 与 $\eps^{kij}=\eps^{jik}$。可见 $e^{jk}$ 是一个反对称的二阶张量，
它把线矢 $m_j$ 变为点矢 $a^k$。限制映射与射影映射都把三维元素变为二维元素，因此都不可逆。

以上讨论的是点列。若把点矢与对偶矢量的地位互换，就得到关于线束的嵌入、限制与射影，结论完全平行，
此处不再重复。

\section{二维矢量空间的特殊性}

二维矢量空间中的矢量、对偶矢量以及上下指标的定义与三维完全相同。差别在于维数少了一，这使
二维空间具有一些在三维中不成立的性质。

\begin{proposition}\label{prop:2d-unique}
设 $n$ 是二维空间上的非零对偶矢量，则方程 $n_ia^i=0$ 的解构成一条过原点的直线，即解空间是
一维的；它的一个（非零）解是
\[
a^i=e^{ij}n_j
\]
（这里用 $e$ 搬运指标，见定义~\ref{def:e}），且一切解都与它成比例。等价地，方程 $a^il_i=0$ 的
解为 $a^i=e^{ij}l_j$，只确定方向而不确定长度。
\end{proposition}

\begin{proof}
取定坐标系，把 $a$ 写成 $a^i=(x,1)$，则 $n_ia^i=0$ 是关于 $x$ 的一次方程
$n_1x+n_2=0$；由于 $n\ne 0$，至少有一个 $n_i$ 非零，故解唯一，解空间至多一维。直接验证：
$n_ie^{ij}n_j=e^{ij}n_in_j=0$，因为 $e^{ij}$ 反对称而 $n_in_j$ 对称。故解空间恰为一维。
\end{proof}

注意与三维的对比：三维中 $n\cdot a=0$ 的解构成一个二维平面，有无穷多个不共线的解；二维中却
只剩下一个方向。

\begin{proposition}\label{prop:2d-conic}
设 $g_{ij}$ 是二维空间上的不定二次型。把 $a$ 写成 $a^i=(x,1)$，则方程
$g_{ij}a^ia^j=0$ 是关于 $x$ 的一元二次方程，故至多有两个解；当 $g$ 不定时恰有两个实根。
即二次型在二维空间中至多确定两个方向。
\end{proposition}

在三维中，$g_{ij}a^ia^j=0$ 确定的是圆锥曲线的点集，一般有无穷多个点。二维中只剩下两个方向，
这一“至多两个解”的事实是后文对合变换只可能有两个不动点的根源。

在二维中 $\eps_{ijk}$ 没有意义：它的非零分量要求三个指标互不相同，而二维指标只取 $1,2$，因此
该张量恒为零。替代它的是反对称二阶张量。

\begin{definition}[二维反对称张量 $e_{ij}$]\label{def:e}
二维空间上非零的反对称二阶张量记为 $e_{ij}$ 与 $e^{ij}$，满足
\[
e_{ij}=-e_{ji},\qquad e^{ij}=-e^{ji},\qquad e_{ij}e^{ij}=2 .
\]
在一定坐标系下取 $e_{12}=1$，则
\[
(e_{ij})=\begin{pmatrix}0&1\\-1&0\end{pmatrix},\qquad
(e^{ij})=\begin{pmatrix}0&1\\-1&0\end{pmatrix},
\]
即有 $e^{12}=1$。此时恒有
\[
e_{ij}e^{ik}=\delta^k_j .
\]
\end{definition}

$e_{ij}$ 是二维空间中最基本的张量之一，它与度规 $g_{ij}$ 的区别在于对称性：$g_{ij}$ 对称，
$e_{ij}$ 反对称。

\subsection{二维叉乘与上下指标的转换}

\begin{definition}[二维叉乘]
两个二维矢量 $a,b$ 的\emph{叉乘}定义为
\[
a\times b=a^ie_{ij}b^j ,
\]
它是一个标量（不是矢量），简记为
\[
[ab]=a^ie_{ij}b^j .
\]
由反对称性立刻得到 $[ab]=-[ba]$，特别地 $[aa]=0$。
\end{definition}

用一个矢量 $a^i$ 与 $e_{ij}$ 缩并，得到一个对偶矢量。记
\[
a_j=a^ie_{ij},\qquad\text{反之}\qquad a^i=e^{ij}a_j .
\]
这与用度规升降指标形式上相似，实质不同：这里用来转换的是反对称张量而不是对称张量，因此
\emph{必须区分 $e_{ij}$ 两个指标的左右}。用几何语言说，$a^i$ 表示二维点矢（点），$a_i$ 表示
二维线矢（过原点的直线）；一个点矢可以画成点，也可以画成它的像直线，二者只是画法不同。

$a^i$ 与 $a_i$ 的这种关系，正是三维中\emph{射影}与\emph{限制}关系的二维版本。给定直线 $l$ 与
$m$：把 $m$ 限制到 $l$ 上得到二维线矢，相当于丢掉 $m$ 上不属于 $l$ 的部分，只留下交点附近的
方向，因此二维线矢宜画成一段直线；把 $l,m$ 作交点（射影）得到二维点矢，因此二维点矢宜画成
一个点。同一个元素的这两种画法由 $e$ 相互转换，纯粹是画法不同而已。

\section{二维伴随矩阵与指标升降}

二维情形下求逆矩阵（更准确地说，求伴随矩阵）有十分简洁的张量表达。设
\[
M=\begin{pmatrix}a&b\\c&d\end{pmatrix},
\]
其伴随矩阵（等于行列式乘以逆矩阵）为
\[
M^*=\begin{pmatrix}d&-b\\-c&a\end{pmatrix},\qquad MM^*=M^*M=(\det M)I .
\]
$M$ 与 $M^*$ 的关系只是把元素换位并改变符号，这说明它可以由 $e_{ij}$ 简单地表达。

\begin{lemma}\label{lem:adj}
设 $M_{ij}$ 是二维空间上的 $(0,2)$ 型张量，$\det M=1$，用 $e$ 升降指标。则 $M$ 与它的伴随
矩阵之间的指标关系为
\[
(M^*)^{ij}=-e^{ik}M_{kl}e^{jl},
\]
并且把 $M^{-1}$ 记作 $M^{ij}$ 时有
\[
M_{ij}M^{ik}=\delta^k_j,\qquad M_{ji}M^{jk}=\delta^k_i .
\]
\end{lemma}

\begin{proof}
\textbf{第一步：矩阵恒等式。}
记 $e=\begin{pmatrix}0&1\\-1&0\end{pmatrix}$，即 $e_{ij}$ 的矩阵，于是 $e^{T}=-e$。取
$M=\begin{pmatrix}a&b\\c&d\end{pmatrix}$，直接计算：
\[
e\,M\,e^{T}
=\begin{pmatrix}0&1\\-1&0\end{pmatrix}\begin{pmatrix}a&b\\c&d\end{pmatrix}
\begin{pmatrix}0&-1\\1&0\end{pmatrix}
=\begin{pmatrix}d&-c\\-b&a\end{pmatrix}
=\operatorname{adj}(M)^{T},
\]
其中 $\operatorname{adj}(M)=M^*=\begin{pmatrix}d&-b\\-c&a\end{pmatrix}$。

\textbf{第二步：写成指标。}
上面等式的 $(i,j)$ 分量为 $(eMe^{T})_{ij}=e_{ik}M_{kl}e_{jl}$，而
$(\operatorname{adj}M)^T_{ij}=(\operatorname{adj}M)_{ji}=(M^*)^{ji}$，故
\[
(M^*)^{ji}=e_{ik}M_{kl}e_{jl}.
\]
把 $i,j$ 都升起来（每升一次指标用 $e$ 搬运，共两次符号变化），并注意 $\det M=1$ 时
$M^{ij}=e^{ia}e^{jb}M_{ab}$ 所表示的矩阵恰为 $(M^{-1})^{ij}$（这正是 $eMe^{T}=\operatorname{adj}(M)^T
=\det M\cdot(M^{-1})^{T}$ 的指标形式），得
\[
(M^*)^{ij}=-e^{ik}M_{kl}e^{jl}.
\]

\textbf{第三步：伴随与逆。}
由 $\det M=1$ 有 $M^*=M^{-1}$，又由 $(M^*)^{ji}=e_{ik}M_{kl}e_{jl}=-(M^*)^{ij}$ 型的指标搬运可得
\[
(M^*)^{ij}=e^{ia}e^{jb}(M^*)_{ab}=e^{ia}(M^{-1})_{ab}e^{jb} .
\]
把 $M^{-1}$ 记为具有两个上指标的 $M^{ab}$，并约定缩并对齐指标位置，即得
\[
M_{ij}M^{ik}=\delta^k_j .
\]
再把 $M_{ij}$ 换成它的转置 $M_{ji}$，由($e$ 的反对称性带来的两次符号变化)同样有
\[
M_{ji}M^{jk}=\delta^k_i .
\]
\end{proof}

\begin{remark}
搬运二维下指标的是反对称的 $e$，不是对称的度规 $g$。因此 $M^{ij}$ 与 $M_{ij}$ 并不互为转置，
一般的张量满足的是上面这类“一对指标上下相消、另一对左右相消”的等式；把 $e$ 与 $g$ 混用、
或随意交换 $M$ 的两个指标都会出错。本节一律用 $e$ 搬运二维指标。
\end{remark}

由引理还可以得到常用的变形。由 ${M^j}_l=M^{jk}e_{kl}$ 与 $M_{ji}M^{jk}=\delta^k_i$ 联立，
\[
M_{ji}{M^j}_l=M_{ji}M^{jk}e_{kl}=\delta^k_ie_{kl}=e_{il},
\]
类似地，把两个指标都换成上指标后有
\[
{M_j}^k{M^j}_l=-\delta^k_l .
\]

\section{射影变换与对合变换}

\subsection{线性映射与射影变换}

在第二章中已经说明，$(1,1)$ 型张量给出矢量到矢量的线性映射。设 $F$ 是二维矢量空间上的线性映射，
其张量写作 ${f^i}_j$，则
\[
F(a)^i={f^i}_ja^j .
\]
本节总要求 $F$ 可逆，即 $\det({f^i}_j)\ne0$。

线性映射的典型例子是点列与线束之间的投影。取定一个不在直线 $L$ 上的点 $A$：对 $L$ 上的动点 $D$，
作直线 $AD$，得到一个过 $A$ 的直线；反过来，对过 $A$ 的动直线 $AD$，作它与 $L$ 的交点 $D$，
得到 $L$ 上的点。这两个对应都把点列 $L$ 与以 $A$ 为中心的线束一一对应，并且都是线性的，因为
“过两点作直线”与“求两直线交点”分别由叉乘的线性运算实现。

线性映射的复合仍是线性映射。例如把 $L$ 上的点 $D$ 经由 $A$ 投影到另一条直线 $L'$ 上的点 $E$，
再把 $E$ 经由另一定点 $A'$ 投影到 $L$ 上的点 $G$，则“$D\mapsto G$”是两次投影的复合。由于
线性映射总可以分解为若干次这样的投影，我们也把可逆线性映射称为\emph{射影变换}。

\begin{definition}[射影变换]
二维矢量空间上可逆的线性映射 $F$，以及由它诱导的点列之间、线束之间的一一对应，统称为
\emph{射影变换}。$F$ 与 $\lambda F$（$\lambda\ne0$）诱导同一个几何对应，因为矢量的长度
（即齐次坐标的规范化）不具有几何意义。
\end{definition}

如果一个射影变换把一条直线（或一个线束的中心所成的集合）映射到自身，则这一变换一般有两个
不动点，这正是命题~\ref{prop:2d-conic} 的应用：求不动点归结为解一个二次方程。

\section{交比}

\subsection{交比的定义}

我们希望用若干二维点矢构造出坐标系无关、规范无关的实数。求坐标系无关，只需通过张量积与缩并
得到标量；求规范无关，只需使每个矢量的次数上下相抵。

使用一个或两个点矢做不到：$[aa]=0$ 恒为零，$[ab]$ 在把 $a$ 换成 $\lambda a$ 时改变 $\lambda$
倍。使用四个点矢则可以构造出
\[
J_{ab,cd}=\frac{[ac][bd]}{[ad][bc]} ,
\]
它关于每个矢量的伸缩都是零次的：$a\to\lambda a$ 使分子分母各出现一次 $\lambda$，因而相消。
这个量称为四点 $a,b,c,d$ 的\emph{交比}，在不引起混淆时简记为 $J$。

\begin{definition}[交比]
设 $a,b,c,d$ 是二维空间中的四个点矢，且
$[ad]\ne0$、$[bc]\ne0$。称
\[
J_{ab,cd}=\frac{[ac][bd]}{[ad][bc]}
\]
为这四点的\emph{交比}。若其中某个叉乘为零，则交比取 $0$ 或 $\infty$。
\end{definition}

若两条直线上各取四个点，且它们由同一个射影变换一一对应，则由于射影变换是线性的而交比由缩并
构成，对应四点的交比相同。

\begin{proposition}\label{prop:J-inv}
设 $F$ 是二维空间上的可逆线性映射，则
\[
J_{F(a)F(b),F(c)F(d)}=J_{ab,cd}.
\]
即交比在射影变换下不变。
\end{proposition}

\begin{proof}
设 $F$ 的张量为 ${f^i}_j$，则 $F(a)^i={f^i}_ja^j$。由行列式的定义，
$[F(a)F(b)]=\det(f)[ab]$，即叉乘在 $F$ 作用下统一乘上因子 $\det f$。代入交比的定义，分子与
分母各含两个因子 $\det f$，故相消。
\end{proof}

\subsection{两个恒等式}

交比的全部代数性质来自下面两个恒等式。

\begin{lemma}[二维叉乘的线性关系]\label{lem:cyc}
对任意三个二维矢量 $a,b,c$ 有
\[
[ab]c+[bc]a+[ca]b=0 .
\]
\end{lemma}

\begin{proof}
三个矢量在二维空间中必线性相关，设 $\alpha a+\beta b+\gamma c=0$ 且系数不全为零。若
$[ab]\ne0$，则 $a,b$ 线性无关，可把它们取作基，把 $c$ 写成 $c=xa+yb$。于是
\[
[bc]=x[ba]=-x[ab],\qquad [ca]=y[ab],
\]
代入左边得
\[
[ab](xa+yb)-x[ab]a+y[ab]b=0 .
\]
若 $[ab]=0$，则 $a,b$ 线性相关，此时 $[bc]a$ 与 $[ca]b$ 都与 $a$ 共线且系数成比例，逐项检查
（或直接用坐标展开）同样得零。
\end{proof}

\begin{lemma}\label{lem:four}
对任意四个二维矢量 $a,b,c,d$ 有
\[
[ab][cd]+[ac][db]+[ad][bc]=0 .
\]
\end{lemma}

\begin{proof}
在引理~\ref{lem:cyc} 中把 $c$ 换成 $d$：
\[
[ab]d+[bd]a+[da]b=0 .
\]
两端与 $c$ 作二维叉乘，即都取 $[\;\cdot\;,c]$。由于 $[uv]$ 对每个自变量都是线性的，左端逐项
展开给出
\[
[ab][dc]+[bd][ac]+[da][bc]=0 .
\]
由反对称性，$[dc]=-[cd]$，$[bd]=-[db]$，$[da]=-[ad]$，代入并整理得
\[
-[ab][cd]-[db][ac]-[ad][bc]=0,
\]
移项即
\[
[ab][cd]+[ac][db]+[ad][bc]=0 .
\]
\end{proof}

\subsection{交比的性质}

由 $[ab]=-[ba]$（即交换两个自变量使叉乘变号）与引理~\ref{lem:four}，可以导出交比的如下性质。

\begin{proposition}[交比在置换下的变化]\label{prop:J-perm}
设 $J=J_{ab,cd}$，且四个点互不相同（即 $[ab],[ac],\dots$ 均不为零）。则
\begin{enumerate}
\item 交换两组之间的元素保持交比不变：$J_{cd,ab}=J_{ba,dc}=J$；
\item 交换同一组内的两个元素使交比变为倒数：$J_{ba,cd}=J_{ab,dc}=\dfrac1J$；
\item 把两组的元素重新搭配，交比与 $1$ 的差互换：
\[
J_{ac,bd}=J_{bd,ac}=1-J .
\]
\end{enumerate}
\end{proposition}

\begin{proof}
选取基使 $[ab]=1$：取 $a=e_2$、$b=e_1$，则 $e_{21}=e_{12}=1$。把 $c,d$ 写成
\[
c=x_ca+y_cb,\qquad d=x_da+y_db .
\]
由第一、三节给的二维叉乘，$[uv]=u^ie_{ij}v^j$ 对两个自变量都是线性的，于是
\[
[ac]=y_c,\quad [bc]=-x_c,\quad [ad]=y_d,\quad [bd]=-x_d,\quad [ab]=1,
\]
以及
\[
[cd]=[x_ca+y_cb,\ x_da+y_db]=x_cy_d[a b]+y_cx_d[b a]=x_cy_d-x_dy_c .
\]
于是
\[
J=J_{ab,cd}=\frac{[ac][bd]}{[ad][bc]}=\frac{y_c(-x_d)}{y_d(-x_c)}=\frac{x_d}{x_c}\cdot\frac{y_c}{y_d}.
\]

第一款：交换两组元素只是把两个点对换个位置，计算得
\[
J_{cd,ab}=\frac{[ca][db]}{[cb][da]}=\frac{(-y_c)(x_d)}{(x_c)(-y_d)}=J,
\]
同理 $J_{ba,dc}=J$。

第二款：由定义 $J_{ba,cd}=\dfrac{[bc][ad]}{[bd][ac]}$，这正是 $J_{ab,cd}$ 的倒数，故
$J_{ba,cd}=1/J$。同理交换 $c,d$ 得 $J_{ab,dc}=\dfrac{[ac][db]}{[ab][cd]}$，它与 $J$ 的比值为
$1/J$，故也等于 $1/J$。（也可直接由第一款推得：把第一款用于以 $b,a$ 为前两点的交比即得此式。）

第三款：由上面的坐标式直接计算
\[
J_{ac,bd}=\frac{[ab][cd]}{[ad][cb]}=\frac{[cd]}{y_d\,x_c}
=\frac{x_cy_d-x_dy_c}{x_cy_d}=1-\frac{x_dy_c}{x_cy_d}=1-J .
\]
同样 $J_{bd,ac}=1-J$。
\end{proof}

\begin{remark}\label{rem:J-not-1}
四个点互不相同时，交比不可能取 $0,1,\infty$。$J=0$ 要求分子 $[ac][bd]=0$，$J=\infty$ 要求
分母 $[ad][bc]=0$，二者都与四点互异（任意两个不成比例，故任意叉乘非零）矛盾。$J=1$ 时由
命题~\ref{prop:J-perm} 第三款得 $J_{ac,bd}=1-J=0$，同样矛盾。因此四点互异时 $J$ 是
$0,1,\infty$ 以外的实数。
\end{remark}

\section{调和分割}

\begin{definition}[调和分割]
设 $J_{ab,cd}=-1$，则称 $a,b$ 与 $c,d$ 互相\emph{调和分割}，称 $a,b,c,d$ 构成\emph{调和点列}。
\end{definition}

由命题~\ref{prop:J-perm}，$J=-1$ 这一条件在所有置换下保持形式（$1/J=-1$，$1-J=2$，而重新
搭配元素得到的 $-1$ 仍落在调和的条件里），因此“哪两点与哪两点调和分割”这一说法的分组方式
是确定的。

\section{对合变换}

\subsection{对合变换与它的不动点}

\begin{definition}[对合变换]
二维矢量空间上的可逆线性映射 $F$ 称为\emph{对合变换}，若存在非零实数 $\lambda$ 使得
\[
F(F(a))=\lambda a\qquad\text{对一切 }a\text{ 成立}.
\]
\end{definition}

对合变换的意义是它把所有的矢量两两配对：$a$ 与 $F(a)$ 配成一对，称为互相对合。若 $F$ 不是
对合变换，则由一对 $a,F(a)$ 可以区分原像与像；若 $F$ 是对合变换，则用 $F$ 作用这一对元素时它们
同时变成对方，原像与像的地位完全对称，因而这一对元素确实被“配对”了。

既然对合变换中原像与像的地位对称，可以期望它有特别对称的代数形式。结论是它可以由一个对称的
$(0,2)$ 型张量决定。

\begin{theorem}[对合变换由对称可逆张量决定]\label{thm:involution}
设 $g_{ij}$ 是二维空间上的对称、可逆 $(0,2)$ 型张量，用 $e$ 升降二维指标。定义线性映射
\[
F(a)^k=g_{ij}a^ie^{jk}\qquad\text{即 }F(a)^k={f^k}_ia^i,\quad {f^k}_i=g_{ij}e^{jk}.
\]
则 $b=F(a)$ 当且仅当 $g_{ij}a^ib^j=0$，并且
\begin{enumerate}
\item $\det F=-\det g\ne0$，故 $F$ 可逆；
\item $F(F(a))=-\det(g)\,a$ 对一切 $a$ 成立，特别地 $F$ 是对合变换。
\end{enumerate}
\end{theorem}

\begin{proof}
\textbf{第一步：$b=F(a)$ 等价于配对条件。}
对任意 $c$ 计算
\[
g_{cd}a^cF(a)^d=g_{cd}a^cg_{ij}a^ie^{jd}.
\]
把哑指标 $c$ 与 $i$ 互换：$a^ca^i$ 对称、$e^{jd}$ 与 $c,i$ 无关，而 $g_{cd}g_{ij}$ 互换 $c,i$
后不变，故右端等于自身的相反数，从而必为零。于是 $g_{ij}a^ib^j=0$。反过来，由
命题~\ref{prop:2d-unique}，方程 $g_{ij}a^ib^j=0$ 关于 $b$ 的解空间是一维的；又 $F(a)=0$ 只在
$a=0$ 时发生（$a\mapsto g_{ij}a^i$ 与 $b\mapsto b^je^{jk}$ 都可逆），故一切解都与 $F(a)$ 成比例。

\textbf{第二步：用矩阵乘法计算 $F^2$。}
取定坐标，把 $g_{ij}$ 写成对称矩阵 $G=\begin{pmatrix}a&b\\b&d\end{pmatrix}$，记
$\Delta=\det G=ad-b^2\ne0$，并把 $e_{ij}$ 与 $e^{ij}$ 都取成矩阵
$J=\begin{pmatrix}0&1\\-1&0\end{pmatrix}$。由于 ${f^k}_i=g_{ij}e^{jk}$ 中缩并的是 $g$ 的第二个
指标与 $e$ 的第一个指标，$F$ 的矩阵恰为
\[
f=GJ=\begin{pmatrix}a&b\\b&d\end{pmatrix}\begin{pmatrix}0&1\\-1&0\end{pmatrix}
=\begin{pmatrix}-b&a\\-d&b\end{pmatrix}.
\]
于是
\[
f^2=\begin{pmatrix}-b&a\\-d&b\end{pmatrix}^2
=\begin{pmatrix}b^2-ad&0\\0&b^2-ad\end{pmatrix}
=-\Delta\,I,
\]
即 $F(F(a))=-\Delta\,a=-\det(g)\,a$。由于 $\det g\ne0$，$F$ 是对合变换。又
$\det f=\det G\cdot\det J=-\Delta\ne0$，故 $F$ 可逆。
\end{proof}

\begin{remark}
使用记号时务必写清用哪个张量搬运指标：二维点列中用 $e_{ij},e^{ij}$，三维中才用 $\eps_{ijk}$；
若改用度规 $g$ 升降指标，一般公式（例如 $M_{ij}M^{jk}\ne\delta^k_i$）都要相应改变。$e$ 是反对称
张量，写这类表达式时上下指标与前后顺序都不能随意调换，本节中以 $b^k=g_{ij}a^ie^{jk}$ 为准。
定理只断言 $F(F(a))$ 与 $a$ 成正比且比例常数非零，其正负与数值随 $e$ 的归一化约定而变。
\end{remark}

\begin{proposition}[不动点]\label{prop:fixed}
设 $F$ 由对称可逆张量 $g_{ij}$ 按定理~\ref{thm:involution} 给出，则 $a$ 是 $F$ 的不动点
（即 $F(a)$ 与 $a$ 成比例）当且仅当
\[
g_{ij}a^ia^j=0 .
\]
若 $g$ 不定，则 $F$ 恰有两个不动点（两个方向），记作 $p,q$。
\end{proposition}

\begin{proof}
设 $F(a)=\mu a$ 且 $a\ne0$。由定理~\ref{thm:involution} 第一步，
$g_{ij}a^iF(a)^j=g_{ij}a^i(\mu a)^j=0$，即
\[
\mu\,g_{ij}a^ia^j=0 .
\]
因 $\mu\ne0$，得 $g_{ij}a^ia^j=0$。

反之设 $g_{ij}a^ia^j=0$ 且 $a\ne0$。仍由定理~\ref{thm:involution} 第一步，取 $b=F(a)$ 时有
$g_{ij}a^ib^j=0$；而取 $b=a$ 时也成立，即 $a$ 自身也满足这个关于 $b$ 的线性方程
$g_{ij}a^i(\;\cdot\;)^j=0$。由命题~\ref{prop:2d-unique}，该方程的解空间是一维的，故 $F(a)$ 与
$a$ 成比例，即 $a$ 是不动点。

解的个数由命题~\ref{prop:2d-conic} 给出：$g$ 不定时二次方程恰有两个实根，即恰有两个不动点
（两个方向）。
\end{proof}

用坐标写出对合的一般形状。按通常规范化取 $a^i=(x_1,1)$、$b^i=(x_2,1)$，则
$g_{ij}a^ib^j=0$ 是关于 $x_1,x_2$ 的双线性方程
\[
Ax_1x_2+B(x_1+x_2)+C=0 ,
\]
其中 $A,B,C$ 由 $g_{ij}$ 的分量给出。取 $a=b$，不动点满足
\[
Ax^2+2Bx+C=0 ,
\]
这正是二次方程有两个根的体现，也与韦达定理相呼应：两根之积与两根之和分别由 $C/A$ 与 $-2B/A$
给出。对一维线束，把线矢坐标写成 $m_i=(k_1,1)$、$n_i=(k_2,1)$，对合同样表示为
$Ak_1k_2+B(k_1+k_2)+C=0$。

\subsection{对合与调和点列}

\begin{theorem}\label{thm:inv-harmonic}
设对合变换 $F$ 的不动点为 $p,q$，且 $a,b$ 是互相对合的一对点（即 $b$ 与 $F(a)$ 成比例）。
则 $p,q,a,b$ 构成调和点列，即
\[
J_{pq,ab}=-1 .
\]
反之，若 $p,q$ 是给定的两个不同方向，则使 $a,b$ 调和分割 $p,q$ 的配对关系是一个对合变换，
其不动点为 $p,q$。
\end{theorem}

\begin{proof}
\textbf{第一步：把 $J_{pq,ab}=-1$ 化为一个对称双线性方程。}
取 $p,q$ 为基，把 $a,b$ 写成
\[
a=x_ap+y_aq,\qquad b=x_bp+y_bq .
\]
由 $[uv]$ 的双线性与反对称性，
\[
[pa]=y_a[pq],\quad [qa]=-x_a[pq],\qquad
[pb]=y_b[pq],\quad [qb]=-x_b[pq].
\]
代入交比的定义得
\[
J_{pq,ab}=\frac{[pa][qb]}{[pb][qa]}
=\frac{y_a\,(-x_b)}{y_b\,(-x_a)}=\frac{x_by_a}{x_ay_b},
\]
（公因子 $[pq]^2$ 已相消）。因此
\[
J_{pq,ab}=-1 \iff x_by_a+x_ay_b=0 .
\]
另一方面，在基 $p,q$ 下直接展开
\[
\bigl(p_iq_j+p_jq_i\bigr)a^ib^j
=\bigl(p_iq_j+p_jq_i\bigr)(x_ap+y_aq)^i(x_bp+y_bq)^j
=[pq]^2\bigl(x_by_a+x_ay_b\bigr),
\]
其中用到 $p_ip^i=q_iq^i=0$、$p_iq^j[p,q]$ 型的配对（这只需在两个基矢上逐项代入验证）。
于是调和条件等价于
\[
g_{ij}a^ib^j=0,\qquad g_{ij}=p_iq_j+p_jq_i=\sym{p}{q}.
\]

\textbf{第二步：这个配对是对合变换。}
解出 $b$：由第二节的命题，$g_{ij}a^ib^j=0$ 作为 $b$ 的方程解空间是一维的，且
\[
b^k=-g_{ij}a^ie^{jk}
\]
是它的一个解。于是配对由线性映射
\[
F(a)^k=g_{ij}e^{jk}a^i
\]
给出（负号只改变配对中的定向）。由定理~\ref{thm:involution}，$F$ 线性、可逆且
$F(F(a))=\lambda a$，故它是对合变换。

\textbf{第三步：不动点。}
设 $a=xe_1+e_2$，计算 $[F(a),a]=0$ 的条件。由
\[
F(a)^k=g_{1j}e^{jk}x+g_{2j}e^{jk}
\]
以及 $[F(a),a]=e_{2k}F(a)^k\cdot x-e_{1k}F(a)^k$，把上式代入并利用 $e^{ik}e_{jk}=\delta^i_j$
的缩并（等价地利用二维恒等式 $e_{ik}e^{jk}=\delta^j_i$）得
\[
[F(a),a]=-\bigl(g_{11}x^2+2g_{12}x+g_{22}\bigr).
\]
而 $g_{ij}a^ia^j=g_{11}x^2+2g_{12}x+g_{22}$，故 $F(a)$ 与 $a$ 平行当且仅当
$g_{ij}a^ia^j=0$。这与命题~\ref{prop:fixed} 一致，且说明 $g$ 不定时恰有两个不动点。

\textbf{反向结论。}
方程 $J_{pq,ab}=-1$ 关于 $a,b$ 显然对称，由前两步它等价于 $g_{ij}a^ib^j=0$，给出 $a$ 与 $b$
之间的一个可逆线性配对，并且 $a$ 的像再取像与 $a$ 成比例，故配对是对合；又 $a=p$ 或 $a=q$ 时
方程自动成立，故 $p,q$ 是它的不动点。
\end{proof}

\subsection{对合的对称积形式}

调和点列的条件可以写成完全对称的形状。由定义 $[pa]=p_ia^i$、$[qa]=q_ia^i$（这里用
$p_i=p^je_{ji}$），条件 $\frac{[pa][qb]}{[pb][qa]}=-1$ 等价于
\[
[pa][qb]+[pb][qa]=0 .
\]
把 $[pa][qb]=p_ia^iq_jb^j$、$[pb][qa]=p_ib^iq_ja^j$ 代入，左边等于
\[
\bigl(p_iq_j+p_jq_i\bigr)a^ib^j=0 .
\]
因此若记
\[
g_{ij}=p_iq_j+p_jq_i ,
\]
则对合由 $g_{ij}a^ib^j=0$ 给出。采用对称积的简写
\[
g=\sym{p}{q},
\]
其中小括号表示对称张量积。于是得到：

\begin{proposition}\label{prop:inv-fixedpts}
对合变换由它的两个不动点完全确定：确定对合的 $(0,2)$ 型张量 $g$ 恰是两个不动点的对偶矢量的
对称积 $g_{ij}=p_iq_j+p_jq_i$。换言之，$g$ 是 $p,q$ 的对偶矢量的对称积构成的二次线。
\end{proposition}

\begin{proof}
设 $p,q$ 是不动点且 $g_{ij}a^ib^j=0$ 定义对合。对任意一对对合点 $a,b$，由定理
\ref{thm:inv-harmonic} 有 $[pa][qb]+[pb][qa]=0$，即 $(p_iq_j+p_jq_i)a^ib^j=0$。因此
$h_{ij}=p_iq_j+p_jq_i$ 与 $g_{ij}$ 给出同一个对合关系；而在二维中由命题~\ref{prop:2d-unique}，
满足同一配对关系的对称 $(0,2)$ 型张量彼此只差一个常数因子，故 $g$ 与 $h$ 成比例，特别地
$h$ 也正是由 $p,q$ 决定的那个对合张量。
\end{proof}

等式 $[pa][qb]+[pb][qa]=0$ 的结构值得注意：不仅原像与像 $a,b$ 之间是对称的，两个不动点 $p,q$
之间也是对称的，而且把 $(ab)$ 与 $(pq)$ 绑定后交换这两组也是对称的。最后这一点说明：如果把对合的
一对 $a,b$ 当作新的不动点，则原来的 $p,q$ 关于这个新的对合恰好互相对合。

以上讨论只涉及直线上的点列。对过定点的线束讨论完全相同，可以类似地定义四条直线的交比、线束上
的对合以及调和线束。

\section{子空间上的应用}

\subsection{退化二次曲线：过 $C$ 的调和线束}

设平面上有一条定直线 $l$，$a,b$ 两点在 $l$ 上；另有两条直线 $m,n$ 分别交 $l$ 于 $p,q$
（即 $p$ 是 $m$ 与 $l$ 的交点，$q$ 是 $n$ 与 $l$ 的交点）。问：在什么条件下 $a,b,p,q$ 构成
调和点列？

用二维矢量的语言，这一条件即
\[
(p_iq_j+p_jq_i)a^ib^j=0 .
\]
目前 $p,q,a,b$ 都是二维元素，需要改写成三维矢量的等式。$a,b$ 是嵌入平面上的点，因此可以原封
不动地当作三维点矢使用；而 $p_i,q_i$ 本质上是 $m_i,n_i$ 在 $l$ 上的限制。既然限制只是限制了
定义域，而 $a,b$ 同在这一定义域（即 $l$ 对应的子空间）中，故
\[
p_ia^i=m_ia^i,\qquad q_ib^j=n_jb^j ,
\]
于是上述条件可以改写为
\[
(m_in_j+m_jn_i)a^ib^j=0 .
\]
这一表达式中所有对象都是三维的了。

设
\[
h_{ij}=m_in_j+m_jn_i=\sym{m}{n},
\]
即 $h$ 是两条直线 $m,n$ 构成的退化二次线。则条件 $h_{ij}a^ib^j=0$ 的含义是：

\begin{proposition}\label{prop:degenerate}
设 $h=\sym{m}{n}$ 由两条直线 $m,n$ 构成。则 $h_{ij}a^ib^j=0$ 当且仅当 $h$ 所代表的两条直线
与直线 $a\times b$ 的两个交点被 $a,b$ 调和分割。
\end{proposition}

\begin{proof}
设直线 $a\times b$ 与 $m,n$ 的交点分别为 $p,q$，即 $p=m\times(a\times b)$、
$q=n\times(a\times b)$。由定义 $p,q$ 都是 $m,n$ 限制到直线 $a\times b$ 上以后得到的点，而由
上面的改写 $h_{ij}a^ib^j=0$ 等价于 $(p_iq_j+p_jq_i)a^ib^j=0$，即 $p,q$ 与 $a,b$ 调和分割。
\end{proof}

从张量的角度看，$(0,2)$ 型张量可以理解为把两个矢量映为实数的二元函数，因此可以像线性函数那样
对它的定义域作限制。于是
\[
(pq)_{ij}=p_iq_j+p_jq_i
\]
正是 $h_{ij}$ 在 $l$ 上的限制：限制后的张量作用在 $a,b$ 上与限制前相同。

\subsection{圆锥曲线的限制给出对合}

退化二次线限制到 $l$ 上给出 $l$ 上的一个对合。自然的问题是：若 $h$ 不退化而为圆锥曲线 $g$，
它在 $l$ 上的限制是否也给出对合？答案是肯定的。

\begin{theorem}[圆锥曲线在直线上的限制]\label{thm:conic-restrict}
设 $g_{ij}$ 是三维空间中的对称 $(0,2)$ 型张量（圆锥曲线），$l$ 是一条直线对应的二维子空间。
则 $g$ 在 $l$ 上的限制 $g'_{ij}$ 仍是对称的 $(0,2)$ 型张量，因而确定 $l$ 上的一个对合变换：
对合的两点 $a,b$ 由
\[
g_{ij}a^ib^j=0
\]
确定。若 $g$ 与 $l$ 有两个交点 $p,q$，则 $p,q$ 是这个对合的两个不动点。若 $g$ 退化为二次线，
结论同样成立。
\end{theorem}

\begin{proof}
限制不改变对称性：$g'_{ij}$ 是由 $g$ 在 $l$ 的基底上取值得到的对称 $(0,2)$ 型张量。若 $l$ 不与
$g$ 相切，则 $g'$ 可逆，定理~\ref{thm:involution} 说明它确定 $l$ 上的一个对合变换，
配对由 $g_{ij}a^ib^j=0$ 给出。

再证交点是不动点。设 $p,q\in l$ 满足 $g_{ij}p^ip^j=g_{ij}q^iq^j=0$。因为 $p,q$ 都在 $l$ 上，
把 $g$ 换成限制 $g'$ 后作用在 $l$ 中的矢量上结果不变，故
\[
g'_{ij}p^ip^j=g'_{ij}q^iq^j=0 .
\]
由命题~\ref{prop:fixed}，这正是对合的不动点的定义，因此 $p,q$ 是对合的两个不动点。

若 $l$ 与 $g$ 相切，则 $g'$ 退化（秩一），此时配对仍由 $g_{ij}a^ib^j=0$ 定义的对合给出，
但两个不动点重合为切点，是上一段当 $p=q$ 时的极限情形。
\end{proof}

把点与线、点矢与线矢的地位互换，得到对偶命题。

\begin{theorem}[对偶形式：圆锥曲线在点上的限制]\label{thm:conic-dual}
设 $g$ 是圆锥曲线，$a$ 是一个点。则 $g$ 在 $a$ 上的限制（作用于过 $a$ 的直线）确定线束上的一个
对合变换：对合的两条直线 $m,n$ 由
\[
g^{ij}a_im_j=0
\]
确定，其中 $g^{ij}=(g^{-1})^{ij}$ 是 $g$ 的逆（注意此处必须用逆而不能用 $g$ 升指标）。若从 $a$
可向 $g$ 引两条切线，其线矢记为 $p,q$，则 $p,q$ 是对合的不动点。若 $g$ 退化为二次点，
结论同样成立。
\end{theorem}

\begin{proof}
圆锥曲线有两种相互对偶的表示：点表示用 $g_{ij}$，线表示用 $g^{ij}=(g^{-1})^{ij}$；线矢 $m$ 的
极点由
\[
A^i=(g^{-1})^{ij}m_j
\]
给出。在过点 $a$ 的线束中，把直线 $m$ 配到它的\emph{对合线}：先作 $m$ 与点 $a$ 确定的那一点
$A=a\times m$（即 $m$ 上关于 $a$ 的对合方向），再作 $A$ 的极线
\[
n_i=g_{ij}A^j=g_{ij}\eps^{jkl}a_km_l .
\]
由 $n_ia^i=g_{ij}\eps^{jkl}a^ia_km_l=0$（$\eps$ 对 $i,k$ 反对称而 $a^ia_k$ 对称），
$n$ 确实也过 $a$，故这一配对是线束到自身的映射。它是线性的、可逆的，并且作用两次回到 $m$
（因为 $A$ 与 $n$ 的对偶关系是对称的），故为对合变换。再对它降指标可得等价的条件
\[
g^{ij}a_im_j=0 ,
\]
其中必须用 $g$ 的逆来搬运指标。

不动点由 $m=n$ 给出，即 $m$ 与其极点 $A$ 满足 $A\times m$ 与 $m$ 成比例，等价于
$(g^{-1})^{ij}m_im_j=0$。这正是从 $a$ 引出的两条切线所满足的方程（切线的极点即切点，
而切点在切线上），因此两条切线是对合的两个不动点。
\end{proof}

\subsection{极点极线与调和点列}

上述结论给出极点极线的一个基本性质。

\begin{theorem}[极线与调和分割]\label{thm:polar}
若关于圆锥曲线 $g$ 有 $a\cdot b=g_{ij}a^ib^j=0$（即 $a$ 在 $b$ 的极线上，或 $b$ 在 $a$ 的
极线上），则 $a,b$ 调和分割直线 $a\times b$ 与 $g$ 的两个交点 $p,q$。
\end{theorem}

\begin{proof}
记 $g_{aa}=g_{ij}a^ia^j$、$g_{ab}=g_{ij}a^ib^j$、$g_{bb}=g_{ij}b^ib^j$，设 $a,b$ 不成比例。
在直线 $a\times b$ 上取基 $a,b$，把 $p,q$ 写成
\[
p=a+s_pb,\qquad q=a+s_qb .
\]
由 $p,q$ 在圆锥曲线上，$s_p,s_q$ 是方程
\[
g_{aa}+2g_{ab}s+g_{bb}s^2=0
\]
的两个根，故（韦达定理）
\[
s_p+s_q=-\frac{2g_{ab}}{g_{bb}},\qquad s_ps_q=\frac{g_{aa}}{g_{bb}} .
\]

先算交比。由二维叉乘的双线性与反对称性，并记 $C=[ab]\ne0$：
\[
[pa]=[a+s_pb,\ a]=s_p[b a]=-s_pC,\qquad
[pb]=[a+s_pb,\ b]=C,
\]
同理 $[qa]=-s_qC$、$[qb]=C$。于是
\[
J_{pq,ab}=\frac{[pa][qb]}{[pb][qa]}
=\frac{(-s_pC)\cdot C}{C\cdot(-s_qC)}=\frac{s_p}{s_q}.
\]
因此
\[
J_{pq,ab}=-1 \iff s_p=-s_q \iff s_p+s_q=0 \iff g_{ab}=0,
\]
即 $g_{ij}a^ib^j=0$。这正是要证的命题。

（也可直接算 $g_{ij}p^iq^j$：由双线性，
\[
g_{ij}p^iq^j=g_{aa}+g_{ab}(s_p+s_q)+g_{bb}s_ps_q
=g_{aa}+g_{bb}\cdot\frac{g_{aa}}{g_{bb}}=0,
\]
其中用到 $g_{ab}=0$ 与 $s_ps_q=g_{aa}/g_{bb}$；这说明 $p,q$ 也互相对合，与交点是不动点、
$a,b$ 与 $p,q$ 调和分割相一致。）
\end{proof}

\subsection{两个附注}

\begin{remark}[不同的圆锥曲线可以给出相同的限制对合]
不同的圆锥曲线在同一个子空间上的限制可能是相同的。最典型的例子是：椭圆 $g$ 限制到直线 $l$ 上
得到的对合，由 $g$ 与 $l$ 的两个交点 $p,q$ 确定；而任何与 $l$ 交于 $p,q$ 的二次曲线都限制出
同一个对合。

设 $p,q$ 两处的切线构成退化二次线
\[
h_{ij}=l^2g_{ij}-l_il_j,\qquad l^2=g^{ij}l_il_j .
\]
从几何上看，$h$ 与 $l$ 的两个交点正是 $p,q$，因此 $h$ 与 $g$ 在 $l$ 上给出相同的对合。从代数
上看，把 $h$ 限制到直线 $l$（即对满足 $l_ix^i=0$ 的矢量作用）时，第二项 $l_il_j$ 的两因子中必
有一个作用在零矢量上，故第二项的限制恒等于零；于是
\[
h'_{ij}=l^2g'_{ij} .
\]
限制相差一个非零常数因子，给出同一个对合关系。
\end{remark}

\begin{remark}[调和线束]
在上面的构图中，不仅 $p,q,a,b$ 构成调和点列，它们与过某定点 $C$（$C$ 不在直线 $a\times b$ 上）
所连的四条直线还构成调和线束。这是因为从点 $C$ 出发的投影是一个射影变换，而交比在射影变换下
不变，所以点列的调和性与线束的调和性同时成立。

若在 $h=\sym{m}{n}$ 之外再引入 $f=\sym{a}{b}$，则这一关系可以更简洁地写成
\[
f:h=0 ,
\]
即两个对称 $(0,2)$ 型张量的二次缩并为零。
\end{remark}

\section{与诱导度规的类比}

上面的计算可以在微分几何的语言中重新表述。把圆锥曲线对应的度规 $g$ 限制在直线 $l$ 上，得到的
$h$ 用几何的话说就是 $g$ 在“超曲面” $l$ 上的\emph{诱导度规}。事实上，三维空间中由线矢 $l_i$
定义的超曲面 $\{x\mid l_ix^i=0\}$，其诱导度规为
\[
h_{ij}=l^2g_{ij}-l_il_j,\qquad l^2=g^{ij}l_il_j ,
\]
这里 $l_i$ 相当于超曲面的法余矢。直线 $l$ 上的点相当于超曲面上的点，$l$ 的地位相当于超曲面的
法余矢。

在这一对应下：
\begin{enumerate}
\item 超曲面的类时、类空、类光分别对应于直线 $l$ 与椭圆相离、相交、相切；
\item 超曲面上的类光方向对应于直线与椭圆的交点；
\item 类光超曲面的诱导度规是退化的，对应于椭圆在相切的直线上限制得到的对合变换是退化的。
\end{enumerate}

结合前面的命题，当 $l$ 与圆锥曲线相切时，$l^2=0$，于是 $h_{ij}=-l_il_j$ 是退化的对称张量，
它给出的对合只有一个不动点（切点），这正是极限情形下两个不动点重合的表现。
